%% SCRAP: hardware/platform-integration/L4RE_DICTIONARY_ALLOCATION
%% SOURCE: docs/working/hardware/platform-integration/L4RE_DICTIONARY_ALLOCATION.adoc
%% STATUS: WORKING
%% FITS: dev-guide/ch-l4re
%% EDITORIAL: lifted — prose rewritten to press voice

\section{Dictionary Allocation on L4Re}

StarForth currently allocates dictionary entries from the host C library. The
\texttt{vm\_create\_word()} routine calls \texttt{malloc()} for each
\texttt{DictEntry}, and \texttt{FORGET} releases them with \texttt{free()}.
This is correct and well-tested on hosted systems, but it ties the dictionary
to a full libc heap --- an awkward dependency in a microkernel context.

\subsection{Why It Matters for L4Re}

The libc-heap approach carries four liabilities under L4Re: it requires a
complete \texttt{malloc} implementation in the microkernel environment; heap
fragmentation makes allocation behavior non-deterministic; dictionary entries
live outside the VM's own memory model in a separate address space; and the VM
cannot be introspected or snapshotted as a single contiguous region.

A VM-based allocator answers each point. It removes the external dependency, a
simple bump allocator behaves deterministically, everything lives inside one
5\,MB VM memory region, the entire VM state can be serialized at once, and the
model aligns naturally with L4Re dataspaces.

\subsection{Proposed Approach: Conditional Compilation}

A compile-time flag selects between the two strategies, leaving the hosted path
untouched:

\begin{lstlisting}[language=C]
#ifdef STARFORTH_VM_DICT_ALLOC
    DictEntry *entry = dict_alloc_from_vm(vm, total);   /* L4Re/embedded */
#else
    DictEntry *entry = (DictEntry *) malloc(total);     /* hosted */
#endif
\end{lstlisting}

In VM-based mode the 5\,MB region is laid out with a 256\,KB dictionary-header
area at the base (growing upward as entries are allocated), then user data and
compiled code (growing upward from \texttt{HERE}), with block storage filling
the remainder. The allocator itself is a cell-aligned bump pointer:

\begin{lstlisting}[language=C]
static DictEntry *dict_alloc_from_vm(VM *vm, size_t total_bytes) {
    size_t align = sizeof(cell_t);
    total_bytes = (total_bytes + align - 1) & ~(align - 1);
    if (vm->dict_headers_used + total_bytes > DICT_HEADERS_MAX) {
        vm->error = 1;
        return NULL;
    }
    DictEntry *entry = (DictEntry *)(vm->memory + vm->dict_headers_used);
    vm->dict_headers_used += total_bytes;
    return entry;
}
\end{lstlisting}

Because allocation is a bump pointer, \texttt{FORGET} becomes a rewind rather
than a loop of frees: it sets \texttt{dict\_headers\_used} back to the target
entry's offset. VM initialization starts \texttt{dict\_headers\_used} at zero
and offsets \texttt{HERE} to \texttt{DICT\_HEADERS\_MAX}. The fast-lookup cache
is unaffected --- its bucket arrays may still use \texttt{malloc} while the
entries themselves live in VM memory.

\subsection{Scope of Change}

The switch touches only three functions plus one new field. The new field is
\texttt{size\_t dict\_headers\_used} in the \texttt{VM} struct. The functions
are \texttt{vm\_create\_word()} (use the VM allocator), the \texttt{FORGET}
implementation (rewind instead of freeing), and \texttt{vm\_init()} (initialize
the offset and \texttt{HERE}). Estimated effort is roughly six hours: four to
implement and two to test.

\subsection{Current Status}

The present decision is to retain \texttt{malloc()} and document the
abstraction point clearly. The hosted implementation is stable and well-tested,
the two approaches show no performance difference, and the abstraction is
simple enough that the L4Re port can flip the switch when needed. The action
items for that port are to define \texttt{STARFORTH\_VM\_DICT\_ALLOC} in the
L4Re build, implement \texttt{dict\_alloc\_from\_vm()}, convert \texttt{FORGET}
to a rewind, run the full test suite, and verify no leaks in the microkernel
environment.

%% TODO(bob): confirm whether the L4Re port has since adopted VM-based allocation; source dated 2025-10-01, implementation deferred.
